% block: matrix-roots-heading page:206 type:section
\subsection{{矩阵平方根与矩阵 $p$ 次根}}\label{sec:matrix-roots}

% block: matrix-square-root-heading page:206 type:subsection
\subsubsection{{矩阵平方根}}\label{subsec:newton-matrix-square-root}

% block: principal-scalar-square-root page:206 type:ordinary-paragraph
{实数 $a>0$ 的平方根为 $a^{1/2}$ 与 $-a^{1/2}$, 其主平方根为 $a^{1/2}$. 主平方根的概念如下推广到复数: 设 $z$ 是复数, 对 $z=re^{i\theta}$ $(r>0,-\pi<\theta<\pi)$, 定义 $z^{1/2}=r^{1/2}e^{i\theta/2}$. 这称为 $z$ 的主平方根.}

% block: infinitely-many-matrix-square-roots page:206 type:ordinary-paragraph
{现在考虑矩阵平方根的情形, 并定义 $A$ 的主平方根. 如果 $X^2=A$, 则称 $X$ 为 $A$ 的一个平方根. 矩阵的平方根与复数或实数的平方根的区别在于, 前者可能有无穷多个. 例如, 设}
\[S_1=\begin{bmatrix}1&0\\0&1\end{bmatrix},\quad S_2=\begin{bmatrix}-1&0\\0&-1\end{bmatrix},\quad S(\theta)=\begin{bmatrix}\cos\theta&\sin\theta\\\sin\theta&-\cos\theta\end{bmatrix}.\]
{则对任意 $\theta$, 都有 $S_1^2=S_2^2=S(\theta)^2=I$. 这些都是单位矩阵的平方根. 因此, 这是一个具有无穷多个平方根的例子. 虽然平方根有无穷多个, 但特征值位于复平面右半平面的平方根是唯一的. 这个唯一的平方根称为矩阵的主平方根. 在上述例子中, $S_1$ 是单位矩阵 $I$ 的主平方根.}

% block: principal-square-root-transition page:206 type:ordinary-paragraph
{下一个定理给出矩阵主平方根的概念.}

% block: principal-matrix-square-root-theorem page:207 type:theorem
\Needspace{14\baselineskip}
\begin{Theorem}{{矩阵主平方根}}\label{thm:principal-matrix-square-root}\label{srcpage:printed-195}
{设 $A$ 是复方阵, 其特征值 $(\lambda_k\ne0)$ 不位于复平面的负实轴上. 则 $A$ 存在唯一一个平方根, 使得该平方根的所有特征值都位于复平面的右半平面.}
\end{Theorem}

% block: principal-square-root-notation page:207 type:ordinary-paragraph
{在 Theorem~\ref{thm:principal-matrix-square-root} 中, 这个唯一的矩阵平方根称为 $A$ 的矩阵主平方根, 记为 $A^{1/2}$.}

% block: square-root-algorithm-transition page:207 type:ordinary-paragraph
{以下介绍计算矩阵 (主) 平方根的数值算法, 包括直接法和迭代法.}

% block: schur-square-root-derivation page:207 type:ordinary-paragraph
{首先介绍一种著名的基于 Schur 分解的直接法. Schur 分解是将矩阵 $A$ 分解为 $QTQ^{\mathrm H}$, 其中 $T$ 是上三角矩阵, $Q$ 是酉矩阵. 如果能够计算 $U:=T^{1/2}$, 则 $(QUQ^{\mathrm H})^2=QU^2Q^{\mathrm H}=QTQ^{\mathrm H}=A$. 因此, 矩阵 (主) 平方根为 $A^{1/2}=QUQ^{\mathrm H}$. 由于 $U$ 也是上三对角矩阵, 可以通过依次求解 $U^2=T$ 得到 $U$. 所得算法见 Algorithm~\ref{alg:schur-matrix-square-root}. 这里, 矩阵 $T$ (及 $U$) 的 $(i,j)$ 元素记为 $t_{ij}$ (及 $u_{ij}$).}

% block: schur-square-root-algorithm page:207 type:ordinary-paragraph
\begin{algorithm}[H]\caption{{(直接法) 求 $X=A^{1/2}$ 的 Schur 方法}}\label{alg:schur-matrix-square-root}\begin{algorithmic}[1]
\State {计算 $A$ 的 Schur 分解, 即 $A=QTQ^{\mathrm H}$.}
\For{$i=1,2,\ldots,n$}
\State $u_{ii}=t_{ii}^{1/2}$
\EndFor
\For{$j=2,3,\ldots,n$}
\For{$i=j-1,j-2,\ldots,1$}
\State $u_{ij}=\frac{1}{u_{ii}+u_{jj}}(t_{ij}-\sum_{k=i+1}^{j-1}u_{ik}u_{kj})$
\EndFor
\EndFor
\State $X=QUQ^{\mathrm H}$
\end{algorithmic}\end{algorithm}

% block: schur-square-root-no-breakdown page:207 type:ordinary-paragraph
{假设矩阵 $A$ 满足 Theorem~\ref{thm:principal-matrix-square-root} 中的条件. 由 Algorithm~\ref{alg:schur-matrix-square-root} 的第 3 行可知, $u_{ii}=t_{ii}^{1/2}$ $(i=1,2,\ldots,n)$, 且 $u_{ii}$ 是复数 $t_{ii}$ 的主平方根. 根据对矩阵 $A$ 的假设, 所有 $u_{ii}$ 都位于复平面的右半平面. 因此, Algorithm~\ref{alg:schur-matrix-square-root} 第 7 行的分母 $u_{ii}+u_{jj}\ne0$. 所以 Algorithm~\ref{alg:schur-matrix-square-root} 不会发生中断.}

% block: newton-square-root-introduction page:207 type:ordinary-paragraph
{下面介绍一些迭代方法. 这些迭代方法可视为求解 $X^2=A$ 的 Newton 方法 (Newton--Raphson 方法). 现在给出矩阵平方根的 Newton 方法的一种著名推导.}

% block: newton-square-root-error-equation pages:207,208 type:ordinary-paragraph
{设 $X_k$ 是 $A^{1/2}$ 的近似解, $E$ 是相应的误差矩阵, 满足 $X_k+E=A^{1/2}$. 如果能够得到 $E$, 使得\newline{}$(X_k+E)^2=(A^{1/2})^2\Leftrightarrow X_k^2+EX_k+X_kE+E^2=A$,\newline{}则有 $X_k+E=A^{1/2}$. 但由于存在附加项 $EX_k+X_kE$, 求解原问题可能变得更加困难. 另一方面, 由假设 $X_k\approx A^{1/2}$, 可以预期 $E\approx O$ (零矩阵). 因此, 不考虑上面的校正方程, 而考虑从 $X_k^2+E_kX_k+X_kE_k=A$ 求出 $E_k$, 这是 \eqref{eq:sylvester-equation} 中的 Sylvester 方程. 于是得到如下 Newton 方法:}
\begin{enumerate}[label=\arabic*.]
\item {设置初始猜测 $X_0$,}
\item {对 $k=1,2,\ldots$, 直到收敛,}
\item {求解 $E_kX_k+X_kE_k=A-X_k^2$ 以得到 $E_k$,}
\item $X_{k+1}=X_k+E_k$.
\item {结束}
\end{enumerate}

% block: newton-square-root-commuting-initial-guess page:208 type:ordinary-paragraph
{如果选择初始猜测 $X_0$ 使得 $AX_0=X_0A$, 则 Sylvester 方程得到简化. 在这一选择下, 可以证明 $E_kX_k=X_kE_k$. 于是由\newline{}$E_k=(1/2)X_k^{-1}(A-X_k^2)=(1/2)(X_k^{-1}A-X_k)$,\newline{}可得 $X_k+E_k=(1/2)(X_k+X_k^{-1}A)$. 现在在 Algorithm~\ref{alg:newton-matrix-square-root} 中给出矩阵平方根的 Newton 方法.}

% block: newton-square-root-frechet-citation page:208 type:ordinary-paragraph
{上述推导的细节以及基于 Fréchet 导数的推导见 \BookCite{96}.}

% block: newton-square-root-algorithm page:208 type:ordinary-paragraph
\begin{algorithm}[H]\caption{{(迭代法) 求 $X=A^{1/2}$ 的 Newton 方法}}\label{alg:newton-matrix-square-root}\begin{algorithmic}[1]
\State {选择 $X_0$ 使得 $AX_0=X_0A$, 例如 $X_0=I$.}
\For{{$k=0,1,\ldots$, 直到收敛}}
\State $X_{k+1}=\frac{1}{2}(X_k+X_k^{-1}A)$
\EndFor
\end{algorithmic}\end{algorithm}

% block: newton-square-root-convergence page:208 type:ordinary-paragraph
{设 $A$ 是满足 Theorem~\ref{thm:principal-matrix-square-root} 中条件的矩阵. 如果 $A^{1/2}X_0$ 的所有特征值都位于复平面的右半平面, 则 $X_k$ 二次收敛到矩阵平方根 $A^{1/2}$ [\hyperref[bib:source-97]{\ref*{bib:source-97}}, p.140].}

% block: newton-square-root-polynomial-initialization page:208 type:ordinary-paragraph
{在 Algorithm~\ref{alg:newton-matrix-square-root} 中, 初始猜测的选择受到限制, 即 $AX_0=X_0A$. 通常选择 $X_0=A$ 或 $X_0=I$. 由于矩阵多项式 $c_0I+c_1A+c_2A^2+\cdots+c_nA^n$ 与 $A$ 可交换, 自然可以选择\newline{}$X_0(=c_0I+c_1A+c_2A^2+\cdots+c_nA^n)$,\newline{}使得残差的范数 $\|X_0^2-A\|$ 最小. 文献 \BookCite{132} 提出了利用 Gröbner 基和低次矩阵多项式进行近似最小化的方法.}

% block: incremental-newton-square-root-transition page:208 type:ordinary-paragraph
{Algorithm~\ref{alg:newton-matrix-square-root} 的一种数值稳定变体是增量 Newton (IN) 迭代 \BookCite{101}, 见 Algorithm~\ref{alg:incremental-newton-square-root}.}

% block: incremental-newton-square-root-algorithm page:208 type:ordinary-paragraph
\begin{algorithm}[H]\caption{{(迭代法) 求 $X=A^{1/2}$ 的 IN 方法}}\label{alg:incremental-newton-square-root}\begin{algorithmic}[1]
\State $X_0=A$; $E_0=\frac{1}{2}(I-A)$
\For{{$k=0,1,\ldots$, 直到收敛}}
\State $X_{k+1}=X_k+E_k$
\State $E_{k+1}=-\frac{1}{2}E_kX_{k+1}^{-1}E_k$
\EndFor
\end{algorithmic}\end{algorithm}

% block: incremental-newton-square-root-cost page:208 type:ordinary-paragraph
{Algorithm~\ref{alg:incremental-newton-square-root} 的一个特点是矩阵 $A$ 不出现在主迭代, 即第 3 行中. 虽然 Algorithm~\ref{alg:incremental-newton-square-root} 比 Algorithm~\ref{alg:newton-matrix-square-root} 在数值上更稳定, 但它需要额外的矩阵与矩阵乘法. 每个迭代步中, Algorithm~\ref{alg:incremental-newton-square-root} 的计算成本约为 Algorithm~\ref{alg:newton-matrix-square-root} 的 7/4 倍.}

% block: matrix-pth-root-heading page:209 type:subsection
\subsubsection{{矩阵 $p$ 次根}}\label{subsec:matrix-pth-root}\label{srcpage:printed-197}

% block: matrix-pth-root-introduction page:209 type:ordinary-paragraph
{本节介绍求解矩阵 $p$ 次根的直接法和迭代法.}

% block: matrix-pth-root-definition-transition page:209 type:ordinary-paragraph
{满足 $X^p=A$ 的矩阵 $X$ 称为矩阵 $p$ 次根. 如 Section~\ref{subsec:newton-matrix-square-root} 所述, 矩阵平方根可能有无穷多个. 其中有由 Theorem~\ref{thm:principal-matrix-square-root} 刻画的矩阵主平方根. 类似地, 矩阵主 $p$ 次根的概念如下给出:}

% block: principal-matrix-pth-root-theorem page:209 type:theorem
\Needspace{14\baselineskip}
\begin{Theorem}{{矩阵主 $p$ 次根}}\label{thm:principal-matrix-pth-root}
{设 $A$ 是复方阵, 其特征值 $(\lambda_k\ne0)$ 不位于复平面的负实轴上. 则 $A$ 存在唯一一个 $p$ 次根, 使得该 $p$ 次根的所有特征值都位于扇形区域 $\{z\in\mathbb{C}:-\pi/p<\arg(z)<\pi/p\}$.}
\end{Theorem}

% block: principal-matrix-pth-root-notation page:209 type:ordinary-paragraph
{这个唯一的矩阵 $p$ 次根称为 $A$ 的矩阵主 $p$ 次根, 记为 $A^{1/p}$. 如果 $p=2$, 则 Theorem~\ref{thm:principal-matrix-pth-root} 与 Theorem~\ref{thm:principal-matrix-square-root} 等价.}

% block: schur-pth-root-derivation page:209 type:ordinary-paragraph
{以下介绍求 $A^{1/p}$ 的直接法. Schur 方法是一种基于 Schur 分解的直接法. 对矩阵 $A$ 进行 Schur 分解 $(A=QTQ^{\mathrm H})$ 后, 有 $A^{1/p}=QT^{1/p}Q^{\mathrm H}$. 为求得 $T^{1/p}$, 设 $U:=T^{1/p}$. 由于 $U$ 是上三对角矩阵, $U^p$ 也是上三对角矩阵. 于是只需求解 $U^p=T$. 所得算法汇总于 Algorithm~\ref{alg:schur-matrix-pth-root}.}

% block: schur-pth-root-algorithm page:209 type:ordinary-paragraph
\begin{algorithm}[H]\caption{{(直接法) 求 $X=A^{1/p}$ 的 Schur 方法}}\label{alg:schur-matrix-pth-root}\begin{algorithmic}[1]
\State {计算 Schur 分解 $A=QTQ^{\mathrm H}$.}
\For{$j=1,2,\ldots,n$}
\State $u_{jj}=t_{jj}^{1/p}$, $v_{jj}^{(1)}=1$, $v_{jj}^{(k+2)}=u_{jj}^{k+1}$ $(k=0,1,\ldots,p-2)$
\For{$i=j-1,j-2,\ldots,1$}
\For{$k=0,1,\ldots,p-2$}
\State $w_{k+2}=\sum_{\ell=i+1}^{j-1}u_{i\ell}v_{\ell j}^{(k+2)}$
\EndFor
\State $u_{ij}=(t_{ij}-\sum_{k=0}^{p-2}v_{ii}^{(p-k-1)}w_{k+2})/(\sum_{k=0}^{p-1}v_{ii}^{(p-k)}v_{jj}^{(k+1)})$
\For{$k=0,1,\ldots,p-2$}
\State $v_{ij}^{(k+2)}=\sum_{\ell=0}^k v_{ii}^{(k-\ell+1)}u_{ij}v_{jj}^{(\ell+1)}+\sum_{\ell=0}^{k-1}v_{ii}^{(k-\ell)}w_{\ell+2}$
\EndFor
\EndFor
\EndFor
\State $X=QUQ^{\mathrm H}$
\end{algorithmic}\end{algorithm}

% block: pth-root-iteration-transition page:209 type:ordinary-paragraph
{至于迭代方法, Algorithms~\ref{alg:newton-matrix-pth-root} 和~\ref{alg:incremental-newton-pth-root} 分别是 Algorithms~\ref{alg:newton-matrix-square-root} 和~\ref{alg:incremental-newton-square-root} 的推广.}

% block: newton-pth-root-algorithm page:210 type:ordinary-paragraph
\begin{algorithm}[H]\caption{{(迭代法) 求 $A^{1/p}$ 的 Newton 方法}}\label{alg:newton-matrix-pth-root}\label{srcpage:printed-198}\begin{algorithmic}[1]
\State {设置 $p$, 并选择 $X_0$ 使得 $AX_0=X_0A$, 例如 $X_0=A$.}
\For{{$k=0,1,\ldots$, 直到收敛}}
\State $X_{k+1}=\frac{1}{p}[(p-1)X_k+X_k^{1-p}A]$
\EndFor
\end{algorithmic}\end{algorithm}

% block: pth-root-stability-and-variants page:210 type:ordinary-paragraph
{Algorithm~\ref{alg:incremental-newton-pth-root} 的计算成本高于 Algorithm~\ref{alg:newton-matrix-pth-root}. 另一方面, 文献 \BookCite{102} 表明 Algorithm~\ref{alg:incremental-newton-pth-root} 具有数值稳定性. Newton 方法的其他有用变体可参见例如 [\hyperref[bib:source-102]{\ref*{bib:source-102}}, Eqs. (3.6), (3.9)]. 文献 \BookCite{187} 给出了 Algorithm~\ref{alg:incremental-newton-pth-root} 的一种具有数值稳定性且计算成本较低的变体.}

% block: incremental-newton-pth-root-algorithm page:210 type:ordinary-paragraph
\begin{algorithm}[H]\caption{{(迭代法) 求 $A^{1/p}$ 的增量 Newton 方法}}\label{alg:incremental-newton-pth-root}\begin{algorithmic}[1]
\State {设置 $p$, $X_0=I$, $E_0=\frac{1}{p}(A-I)$.}
\For{{$k=0,1,\ldots$, 直到收敛}}
\State $X_{k+1}=X_k+E_k$, $F_k=X_kX_{k+1}^{-1}$
\State $E_{k+1}=-\frac{1}{p}E_k\left[X_{k+1}^{-1}I+2X_{k+1}^{-1}F_k+\cdots+(p-1)X_{k+1}^{-1}F_k^{p-2}\right]E_k$
\EndFor
\end{algorithmic}\end{algorithm}
